\section{The Nine Orders}\label{app:orders}

Scripture names nine orders of angels. Dionysius gathers them from
``the appellations'' of the sacred writers and arranges them in three
triads (\emph{CH} 6); Gregory the Great, preaching on the woman with
the ten drachmas, reads the same nine names in the coin that was found:
\latin{Novem vero angelorum ordines diximus, quia videlicet esse,
testante sacro eloquio, scimus angelos, archangelos, virtutes,
potestates, principatus, dominationes, thronos, cherubim, atque
seraphim} (Gregory, \emph{Hom.\ in Ev.} 34.7). Aquinas sets the two
authorities side by side and holds both orderings reasonable, differing
only in the middle of the scale: Dionysius places the Virtues beneath
the Dominations and the Principalities beneath the Powers; Gregory
places the Principalities beneath the Dominations and the Virtues
beneath the Powers (\ST{I q.~108 a.~6 co.}). Each arrangement can claim
the words of the Apostle: Ephesians 1:21, enumerating upward, supports
Dionysius; Colossians 1:16, enumerating downward, supports Gregory
(\ST{I q.~108 a.~6 co.}). The Dionysian order governs the entries
below (\S\ref{sec:ch}; on the comparative table,
Appendix~\ref{app:orderings}).

\begin{center}
\footnotesize
\begin{tabular}{@{}cll@{}}
\toprule
& \emph{Dionysius} & \emph{Gregory} \\
\midrule
1 & Seraphim & Seraphim \\
2 & Cherubim & Cherubim \\
3 & Thrones & Thrones \\
4 & Dominations & Dominations \\
5 & Virtues & Principalities \\
6 & Powers & Powers \\
7 & Principalities & Virtues \\
8 & Archangels & Archangels \\
9 & Angels & Angels \\
\bottomrule
\end{tabular}
\end{center}

For each order below: the names in Latin, Greek (transliterated), and
Hebrew where a Hebrew name exists; the places of Scripture where the
order is named, quoted from the Douay-Rheims; the characterization in
Dionysius (\emph{CH} 7--9), Gregory (\emph{Hom.\ in Ev.} 34), and
Aquinas (\ST{I q.~108 aa.~5--6}).

\paragraph{1. Seraphim.} \emph{Names:} \latin{seraphim}; \gk{seraphim};
Hebrew \emph{seraphim}, ``burning ones.'' \emph{Scripture:} ``Upon it
stood the seraphims: the one had six wings'' (Isa 6:2); ``one of the
seraphims flew to me, and in his hand was a live coal'' (Isa 6:6);
compare the six-winged living creatures who ``rested not day and night,
saying: Holy, Holy, Holy'' (Apoc 4:8). \emph{Dionysius:} the name
``denotes either that they are kindling or burning,'' their ceaseless
movement about divine things, and their purifying fire (\emph{CH} 7.1).
\emph{Gregory:} \latin{Seraphim namque ardentes vel incendentes
vocantur}; they burn with incomparable love from their singular
nearness to their Maker (34.10). \emph{Aquinas:} the name comes not
from charity only but from its excess, expressed by ardor or fire
(\ST{I q.~108 a.~5 ad~5}).

\paragraph{2. Cherubim.} \emph{Names:} \latin{cherubim};
\gk{cheroubim}; Hebrew \emph{keruvim}. \emph{Scripture:} God ``placed
before the paradise of pleasure Cherubims, and a flaming sword'' (Gen
3:24); he ``sittest upon the cherubims'' (Ps 79:2; 4~Kgs 19:15); the
living creatures of the vision ``are cherubims'' (Ezek 10:20).
\emph{Dionysius:} the name denotes ``a fulness of knowledge or stream
of wisdom,'' readiness to receive the highest gift of light, and its
ungrudging communication (\emph{CH} 7.1). \emph{Gregory:}
\latin{Cherubim quoque plenitudo scientiae dicitur} (34.10).
\emph{Aquinas:} the Cherubim have the excellence of knowledge, the
Seraphim of ardor (\ST{I q.~108 a.~5 ad~6}).

\paragraph{3. Thrones.} \emph{Names:} \latin{throni}; \gk{thronoi};
no Hebrew order-name (Latin \latin{sedes}, ``seats''). \emph{Scripture:}
``whether thrones, or dominations, or principalities, or powers'' (Col
1:16); ``thou hast sat on the throne, who judgest justice'' (Ps 9:5);
the sapphire throne above the cherubim (Ezek 10:1). \emph{Dionysius:}
the name denotes exaltation above every inferiority, an unswerving
settlement around the Highest, and ``their bearing God'' (\emph{CH}
7.1). \emph{Gregory:} thrones are those through whom God exercises
judgment, \latin{quibus ad exercendum judicium semper Deus omnipotens
praesidet} (34.10). \emph{Aquinas:} they excel as having an immediate
knowledge of the types of the divine works; Dionysius explains the name
from material seats (\ST{I q.~108 a.~5 ad~6}).

\paragraph{4. Dominations.} \emph{Names:} \latin{dominationes};
\gk{kyriotetes}; Parker, ``Lordships.'' \emph{Scripture:} Col 1:16;
``above all principality and power and virtue and dominion'' (Eph
1:21). \emph{Dionysius:} the name denotes ``a certain unslavish
elevation,'' a liberty above every servitude, moulding itself and those
after it to the true Lordship (\emph{CH} 8.1). \emph{Gregory:}
\latin{Dominationes autem vocantur qui etiam potestates principatum
dissimilitudine alta transcendunt} (34.10). \emph{Aquinas:} the name
means ``a certain liberty, free from servile condition and common
subjection'' (\ST{I q.~108 a.~5 ad~2}); because the Dominations are
above all subjection, they are not sent to outward ministry (\ST{I
q.~112 a.~4 s.c.}).

\paragraph{5. Virtues.} \emph{Names:} \latin{virtutes}; \gk{dynameis};
Parker, ``Powers.'' \emph{Scripture:} Eph 1:21; ``the angels and powers
and virtues being made subject to him'' (1~Pet 3:22). \emph{Dionysius:}
the name denotes ``a certain courageous and unflinching virility'' for
all the godlike energies, unflinchingly looking to the power-making
power (\emph{CH} 8.1). \emph{Gregory:} \latin{Virtutes etenim vocantur illi nimirum spiritus,
per quos signa et miracula frequentius fiunt} (34.10). \emph{Aquinas:} the name ``signifies a
certain virile and immovable strength,'' first in regard of the divine
operations (\ST{I q.~108 a.~5 ad~1}).

\paragraph{6. Powers.} \emph{Names:} \latin{potestates};
\gk{exousiai}; Parker, ``Authorities.'' \emph{Scripture:} Col 1:16;
``nor angels, nor principalities, nor powers'' (Rom 8:38); ``the
principalities and powers in heavenly places'' (Eph 3:10).
\emph{Dionysius:} the name denotes ``the beautiful and unconfused good
order'' in the divine receptions, an authority used not imperiously but
benignly (\emph{CH} 8.1). \emph{Gregory:} those to whose power the
adverse virtues are subject and restrained, \latin{quorum potestate
refrenantur, ne corda hominum tantum tentare praevaleant quantum
volunt} (34.10). \emph{Aquinas:} the name ``signifies a kind of
ordination'' as regards the reception of divine things and their
execution (\ST{I q.~108 a.~5 ad~3}).

\paragraph{7. Principalities.} \emph{Names:} \latin{principatus};
\gk{archai}. \emph{Scripture:} Col 1:16; Eph 1:21; Eph 3:10; Rom 8:38.
\emph{Dionysius:} the name manifests ``their princely and leading
function, after the Divine example,'' their being wholly turned to the
super-princely Prince (\emph{CH} 9.1). \emph{Gregory:} those who
preside over the good spirits themselves, \latin{qui ipsis quoque bonis
angelorum spiritibus praesunt} (34.10). \emph{Aquinas:} the name
signifies ``one who leads in a sacred order'' (\ST{I q.~108 a.~5
ad~3}).

\paragraph{8. Archangels.} \emph{Names:} \latin{archangeli};
\gk{archangeloi}. \emph{Scripture:} ``with the voice of an archangel
and with the trumpet of God'' (1~Thess 4:15); ``when Michael the
archangel, disputing with the devil'' (Jude 9); Gabriel, ``who stand
before God'' (Luke 1:19); Raphael, ``one of the seven, who stand before
the Lord'' (Tob 12:15). \emph{Dionysius:} the middle order of the last
triad, ``of the messenger Order,'' receiving the divine illuminations
from the first powers and announcing them to the angels (\emph{CH}
9.2). \emph{Gregory:} \latin{qui vero summa annuntiant, archangeli,
vocantur} (34.8). \emph{Aquinas:} the Archangels ``are called the
angel princes; forasmuch as they are princes as regards the `Angels,'
and angels as regards the Principalities'' (\ST{I q.~108 a.~5 ad~4}).

\paragraph{9. Angels.} \emph{Names:} \latin{angeli}; \gk{angeloi};
Hebrew \emph{mal'akhim}, ``messengers.'' \emph{Scripture:} ``their
angels in heaven always see the face of my Father'' (Matt 18:10); the
angels minister to him (Matt 4:11); ``thousands of thousands ministered
to him'' (Dan 7:10); ``Michael and his angels fought with the dragon''
(Apoc 12:7). \emph{Dionysius:} they are ``more properly named Angels by
us than those of higher degree, because their Hierarchy is occupied
with the more manifest, and is more particularly concerned with the
things of the world'' (\emph{CH} 9.2); from them come the guardians of
men (\emph{CH} 5, 9). \emph{Gregory:} \latin{Hi autem qui minima
nuntiant, angeli \dots\ vocantur}; and \latin{angelorum vocabulum,
nomen est officii, non naturae} (34.8). \emph{Aquinas:} the common
name remains proper to the lowest order, since it possesses no
excellence above the common manifestation (\ST{I q.~108 a.~5 ad~1});
the guardian angels are drawn from this order (\ST{I q.~113 a.~3
s.c.}).

Gregory reads the same ninefold company in the figures of the Gospel:
the ninety-nine sheep left in the desert are the highest choirs left in
heaven when the Lord sought the one lost sheep on earth, and the
woman's ten drachmas are nine because the orders of angels are nine
(\emph{Hom.\ in Ev.} 34.3--6). The number of the angels within the
orders exceeds our counting, ``thousand thousands, and myriad myriads''
(Dan 7:10; \emph{CH} 14; \ST{I q.~50 a.~3}).
